\section{Actual moving roots through a growing prime window}
\label{rw-section}

We average over an actual interval of integer roots and retain every
endpoint error in the prime-power congruence counts. Exact rank
stratification and the spacing of its possible primes improve the
window substantially. The conclusion concerns most roots in that
interval; the primes beyond the window and its exceptional roots
remain open.

Write $\mathcal O=\mathbb Z[\zeta]$, $\zeta^2-\zeta+1=0$,
$N(a+b\zeta)=a^2+ab+b^2$ and $P(a+b\zeta)=ab(a+b)$.
For integers $n,B\ge1$ and $B\le k<2B$, put
\[
 w_k=3k+\zeta,\quad Q_k=N(w_k),\quad
 T_{k,n}=|P(w_k^n)|,\quad t_{k,n}=\log c_{k,n},
\]
where $c_{k,n}$ is the largest side of the positive unit rotation
of $w_k^n$. For real $5\le Y\le Z$, define
\[
 E_{k,n}(Y,Z)=\sum_{Y<p\le Z}(v_p(T_{k,n})-3)_+\log p.
\]
This is the primewise positive excess of the actual integer boundary
in the specified finite window.

\begin{lemma}[Actual heights and nonzero boundaries]\label{rw-height}
All these roots and their powers are primitive and unramified,
their boundaries are nonzero, and
\[
 t_{k,n}\ge n\log(3B),\qquad
 \log|P(w_k^m)|\le3m\log(6B+1)\quad(m\ge1).
\]
\end{lemma}
\begin{proof}
The root is primitive and $Q_k=9k^2+3k+1\equiv1\pmod3$.
Oriented Eisenstein factorization preserves primitivity under all
powers. A zero boundary would make $w_k/\bar w_k$ a root of unity;
primitivity and the absence of a ramified factor would then force
$w_k$ to be a unit, contrary to $Q_k\ge13$. Unit rotation preserves
norm and absolute boundary. Its positive coordinates give
$Q_k^n\le c_{k,n}^2$, while $Q_k\ge(3B)^2$.
Also $Q_k\le(6B+1)^2$ and the cubic difference identity yields
$|P(w_k^m)|\le Q_k^{3m/2}$.
\end{proof}

\begin{lemma}[Simple roots after removing the exponent's prime part]
\label{rw-simple}
For $p>3$ and $p\nmid m$, every root modulo $p$ of
$G_m(X)=P((3X+\zeta)^m)\in\mathbb Z[X]$ is simple. There are at
most $3m$ such roots, each lifting uniquely at every precision.
If $a_p=v_p(n)$ and $m_p=n/p^{a_p}$, then at each actual $k$
\[
 v_p(T_{k,n})=
 \begin{cases}
 v_p(|G_{m_p}(k)|)+a_p,&v_p(|G_{m_p}(k)|)>0,\\
 0,&v_p(|G_{m_p}(k)|)=0.
 \end{cases}
\]
\end{lemma}
\begin{proof}
Over the split algebra or the quadratic extension of $\mathbb F_p$,
let $z,\bar z$ be the distinct roots of $X^2-X+1$. The identity
\[
 3(z-\bar z)G_m(X)
   =(3X+z)^{3m}-(3X+\bar z)^{3m}
\]
shows that neither factor $3X+z$, $3X+\bar z$ can vanish at a root.
Their ratio has derivative $3(\bar z-z)/(3X+\bar z)^2\ne0$.
Since $p\nmid3m$, differentiating its $3m$-th power at the root
gives a nonzero result. Thus the root is simple, and the polynomial
degree bound and simple-root lifting prove the first assertions.

If $p\nmid Q_k$, its homogeneous rank is prime to $p$, so it
divides $n$ if and only if it divides $m_p$. The established rank
and lifting law gives the displayed identity. If $p\mid Q_k$,
primitivity implies that $p$ divides neither boundary; both
valuations are zero. Consequently, at every actual root,
\[
 (v_p(T_{k,n})-3)_+
 \le (v_p(|G_{m_p}(k)|)-3)_++v_p(n).
\]
The total additional lifting cost in any window is at most $\log n$.
\end{proof}

\begin{proposition}[An unstratified actual interval estimate]
\label{rw-unstratified}
One has
\[
 \frac1B\sum_{B\le k<2B}\frac{E_{k,n}(Y,Z)}{t_{k,n}}
 \le \frac{10\log Y}{Y^3\log(3B)}
       +\frac{18n\pi(Z)}B+\frac{\log n}{n\log(3B)}.
\]
\end{proposition}
\begin{proof}
For $m=m_p$, each of the at most $r_{p,m}\le3m$ lifted residue
classes meets the interval in at most $B/p^e+1$ integers. The cap
$H=\lfloor3m\log(6B+1)/\log p\rfloor$ covers every actual depth.
Summing the positive-excess layers $4\le e\le H$ gives
\[
 \sum_k(v_p(|G_m(k)|)-3)_+\log p
 \le\frac{3mB\log p}{p^3(p-1)}+9m^2\log(6B+1).
\]
Empty layer ranges contribute zero. Use $m\le n$, divide by
$Bn\log(3B)$, and use $\log(6B+1)\le2\log(3B)$.
The elementary estimate
\begin{equation}\label{rw-prime-series}
 3\sum_{p>Y}\frac{\log p}{p^3(p-1)}
       \le\frac{10\log Y}{Y^3}\qquad(Y\ge5)
\end{equation}
follows by $p/(p-1)\le5/4$ and comparison of the decreasing
function $\log x/x^4$ with its integral from $\lfloor Y\rfloor$.
Its integral is $(\log\lfloor Y\rfloor/3+1/9)/\lfloor Y\rfloor^3$;
$\lfloor Y\rfloor\ge4Y/5$ gives the stated loose constant.
Add the lifting cost from Lemma~\ref{rw-simple}. Every denominator
comparison has a nonnegative numerator.
\end{proof}

\begin{lemma}[The exact rank count]\label{rw-exact-rank}
For $p>3$, put $\chi_p=1$ at split primes and $-1$ at inert primes,
and $N_p=p-\chi_p$. For each integer $d\ge1$, the number of residues $k\bmod p$ for which
$((3k+\zeta)/(3k+\bar\zeta))^3$ has exact order $d$ is
\[
 r_p(d)=
 \begin{cases}
 3\varphi(d)-\mathbf1_{d=1},&d\mid N_p/3,\\
 0,&\text{otherwise}.
 \end{cases}
\]
Only nonzero-norm residues enter this definition. Each counted root
of $G_d$ lifts uniquely to every $p^e$, retaining its exact rank
modulo $p$.
\end{lemma}
\begin{proof}
The fractional-linear map $a\mapsto(a+\zeta)/(a+\bar\zeta)$ is
a bijection onto the norm-one group minus its identity. Its inverse
is $(R\bar\zeta-\zeta)/(1-R)$; conjugation fixes this inverse in
the inert case, and the split case follows directly in the split
algebra. The norm-one group is cyclic of order $N_p$, divisible by
three. Cubing has kernel three and image of order $N_p/3$.
Each order-$d$ element has three preimages, and only at $d=1$ must
the omitted identity be subtracted. Multiplication by three permutes
the parameter residues. Since $p\nmid d$, Lemma~\ref{rw-simple}
applies. Retaining the initial residue retains its rank.
\end{proof}

\begin{theorem}[A rank-stratified actual prime window]
\label{rw-main}
For all the above data,
\begin{equation}\label{rw-final-bound}
 \frac1B\sum_{B\le k<2B}\frac{E_{k,n}(Y,Z)}{t_{k,n}}
 \le \frac{10\log Y}{Y^3\log(3B)}
       +\frac{12(Z+1)}B+\frac{\log n}{n\log(3B)}.
\end{equation}
\end{theorem}
\begin{proof}
At each supported prime use its actual rank $d$ and first depth
$s=v_p(T_{k,d})$. Then $d\mid n$, $d\mid N_p/3$, and
$v_p(T_{k,n})=s+v_p(n)$. It suffices to count $(s-3)_+$, with
at most $\log n$ additional lifting cost at each root.

For each fixed rank-prime pair, Lemma~\ref{rw-exact-rank} gives
at most $3\varphi(d)(B/p^e+1)$ parameters at depth at least $e$.
The cap $\log|G_d(k)|\le3d\log(6B+1)$ yields the bound
\[
 \frac{3\varphi(d)B\log p}{p^3(p-1)}
       +9\varphi(d)d\log(6B+1)
\]
for their first-depth positive cost. At fixed $p$, summing the first
term over $d\mid\gcd(n,N_p/3)$ uses
$\sum\varphi(d)=\gcd(n,N_p/3)\le n$ and
\eqref{rw-prime-series}.

For the endpoint errors, reverse the order of summation. A prime
of rank $d$ lies in $p=3dj+1$ or $p=3dj-1$, with $j\ge1$.
Thus at most $2(Z+1)/(3d)$ candidates occur up to $Z$. The excluded
integer $p=1$ is not counted, including when $d=1$. Therefore
\[
 \sum_{d\mid n}9\varphi(d)d\log(6B+1)\frac{2(Z+1)}{3d}
       =6n(Z+1)\log(6B+1).
\]
Here the depth-cap factor $d$ cancels the progression spacing.
Divide by $Bn\log(3B)$, use $\log(6B+1)\le2\log(3B)$, and add
the lifting cost. Norm-dividing primes have zero boundary valuation,
so the rank partition omits no supported prime.
\end{proof}

\begin{corollary}[A window entering the maximal-rank range]
\label{rw-polynomial}
Take $B=n^4$, $Y=n^{1/6}$ and $Z=n^3$. For sufficiently large $n$,
\[
 \frac1B\sum_k\frac{E_{k,n}(Y,Z)}{t_{k,n}}
 \le\frac5{12\sqrt n}+\frac{12}n+\frac{12}{n^4}+\frac1{4n}.
\]
The proportion for which $E_{k,n}/t_{k,n}>n^{-1/4}$ is at most
\[
 \frac5{12n^{1/4}}+\frac{12+1/4}{n^{3/4}}+\frac{12}{n^{15/4}}.
\]
\end{corollary}
\begin{proof}
Substitute into \eqref{rw-final-bound}, use
$\log(3B)\ge4\log n$, and apply Markov's inequality to the
nonnegative actual cost.
\end{proof}

All these profiles are actual homogeneous profiles with unit residual
and $\lambda=1/n$; their root norms are comparable to $n^8$. The
window includes possible maximal-rank primes even when $n$ is prime.
More generally, $Z=o(B)$ makes the discrepancy vanish whenever the
main and lifting terms vanish. For $B=n^4$ one may take
$Z=\lfloor B/(\log n)^2\rfloor$, obtaining average
$O((\log n)^{-2}+n^{-1/2})$ and exceptional proportion $O(1/\log n)$
at threshold $1/\log n$.

These are finite-window statements about most actual roots. They do
not bound a specified exceptional root, primes above $Z$, or the
complete signed tail of any prescribed orbit. The original
unstratified bound also gives, for every fixed integer $A\ge2$,
vanishing average in $B=n^{A+2}$ through $Z=n^A$. The rank
stratification is stronger and uses no prime-distribution estimate.

The full ordinary proof was independently reviewed, including the
split and inert rank counts, the prime-part removal, endpoint caps,
and the exceptional-proportion constants. Separate exact replays
check actual integer blocks and complete finite residue spaces at
specified precisions. They are diagnostics of the ordinary proof,
not a complete-tail certificate or a Lean formalization.
